\documentclass{report}
\usepackage{amsmath,amssymb}
\setlength{\parindent}{0mm}
\setlength{\parskip}{1em}
\begin{document}
\begin{center}
\rule{6in}{1pt} \
{\large Kirk E. Jordan \\
{\bf A New Day, A New High Performance Computer, What Does It Mean for Iterative Methods?}}

IBM Research \\ 1 Rogers Street \\ Cambridge \\ MA 02142
\\
{\tt kjordan@us.ibm.com}\\
Hormozd Gahvari\\
William D. Gropp \\
Vipin Sachdeva\\
Martin Schulz\\
Ulrike M. Yang\end{center}

A new day is dawning with IBM’s latest High Performance Computing (HPC)
System as we see core counts continue to rise. In this talk, we will
briefly describe this system (in more detail than last year) including
some of the software features and explain how it fits into IBM’s road
to Exascale. More importantly, we will describe some of the work underway
pertinent to Iterative Methods. This will serve as a foundation for a
talk on “Algebraic Multigrid Solvers on Future Supercomputers” by
Hormozd Gahvari. We will describe some of our experience on this new
system including some of our early experiences in using Hybrid Coding
(MPI plus OpenMP). In addition, while related to our HPC work, we will
mention work continuing to make HPC accessible to a wider audience and
eventually targeting the scientific and engineering practitioner directly
and explain why this work is of importance and may be of interest to the
Iterative Methods community.


\end{document}
